% Generated draft: requires mathematical and source review.
\documentclass[11pt]{article}
\usepackage[margin=25mm]{geometry}
\usepackage{fontspec}
\setmainfont{Latin Modern Roman}
\usepackage{amsmath,amssymb,mathtools,unicode-math}
\setmathfont{Latin Modern Math}
\usepackage{enumitem,needspace,xurl,hyperref}
\setlength{\parindent}{0pt}
\setlength{\parskip}{5pt}
\setlength{\emergencystretch}{3em}
\newcommand{\N}{\mathbb{N}}
\newcommand{\Z}{\mathbb{Z}}
\newcommand{\Q}{\mathbb{Q}}
\newcommand{\R}{\mathbb{R}}
\newcommand{\C}{\mathbb{C}}
\newcommand{\F}{\mathbb{F}}
\newcommand{\A}{\mathbb{A}}
\newcommand{\PP}{\mathbb{P}}
\newcommand{\E}{\mathbb{E}}
\newcommand{\eps}{\varepsilon}
\newcommand{\dd}{\,\mathrm{d}}
\begin{document}
\section*{UnsolvedMath ID: 800721}
\textbf{Source problem code:} catalogue entry
\subsection{Definitions and mathematical statement}
Let \(F\) be a field satisfying at least one of the following conditions: \(F\) is finitely generated as a field extension of its prime subfield; or there exists an algebraically closed subfield \(k\subseteq F\) such that \(F/k\) is finitely generated as a field extension. A quaternion division algebra over \(F\) means a central division \(F\)-algebra of dimension \(4\). A finite field extension \(E/F\) splits such an algebra \(D\) if
\[
D\otimes_F E\cong M_2(E)
\]
as \(E\)-algebras. A maximal subfield of \(D\) is a commutative subfield containing \(F\) that is maximal under inclusion; for a quaternion division algebra it has degree \(2\) over \(F\). Two algebras have the same maximal subfields if their maximal subfields have the same \(F\)-isomorphism classes, equivalently if every quadratic field extension \(L/F\) admits an \(F\)-algebra embedding into one precisely when it admits such an embedding into the other.

\textbf{Requested assertion (Problem 35).} For every such field \(F\) and every pair of quaternion division \(F\)-algebras \(D_1,D_2\), prove that
\[
\left[\text{for every finite field extension }E/F,\quad
D_1\otimes_F E\cong M_2(E)
\ \Longleftrightarrow\ 
D_2\otimes_F E\cong M_2(E)\right]
\ \Longrightarrow\ D_1\cong_F D_2.
\]
The source also requests the stronger formulation in which the hypothesis is replaced by equality of maximal subfields:
\[
\left[\text{for every quadratic field extension }L/F,\quad
L\hookrightarrow_F D_1\ \Longleftrightarrow\ L\hookrightarrow_F D_2\right]
\ \Longrightarrow\ D_1\cong_F D_2.
\]
Here \(\hookrightarrow_F\) denotes the existence of an \(F\)-algebra embedding. The supplied statement imposes no characteristic restriction and does not specify separability. The displayed formulations therefore read finite extensions and maximal subfields literally, including inseparable ones when they occur. Whether the original source intended a separability restriction, especially in characteristic \(2\), remains unresolved from the supplied record.
\subsection{Short English statement}
For a field finitely generated over its prime field or over an algebraically closed subfield, the problem asks whether a quaternion division algebra is determined up to isomorphism over that field by its finite splitting fields, or even by its maximal subfields. The supplied record labels the problem open; this draft does not independently verify that status. No characteristic or separability restriction is supplied.
\subsection{Sources}
\begin{itemize}
\item Local UnsolvedMath registry, permanent ID \(800721\), exact source problem code \texttt{catalogue entry}, title \emph{Open Problems on Central Simple Algebras --- Problem 35}, in the AMR Open Problem Lists collection. The record attributes the problem to Auel, Brussel, Garibaldi, and Vishne, with proposed year \(2010\), and supplies the status label \texttt{open}.
\item Auel, Brussel, Garibaldi, and Vishne, \emph{Open Problems on Central Simple Algebras} (2010), Problem 35. Supplied source URL: \url{https://arxiv.org/abs/1006.3304}. The background records access on \(2026\text{-}10\text{-}16\) and extraction by \texttt{arxiv\_latex}. These are supplied provenance statements, not a claim of independent source access or verification.
\item Supplied registry revision: \texttt{d371ee917ac9647fbb06cd6fc529deec6f3dda82}; record updated \(2026\text{-}10\text{-}16\). This timestamp is registry provenance and does not establish a literature assessment as of that date.
\end{itemize}
\end{document}
